\section{Conclusiones}
\label{sec:conclusiones}

La segunda entrega llevó el Sistema de Priorización de Atención de una
implementación inicial a un sistema funcional que cubre todos sus
requerimientos. La cola de prioridad por niveles se conservó, porque las
mediciones confirman que con cinco niveles fijos supera tanto a una lista
ordenada como a un \emph{heap} binario. A ella se agregaron un índice hash
propio y nodos doblemente enlazados que funcionan como localizadores, con
lo que buscar, cancelar y verificar duplicados pasaron de $O(n)$ a $O(1)$
esperado; y la pila del historial ahora permite deshacer la atención más
reciente en $O(1)$, devolviendo la solicitud a su posición original.

La corrección del sistema se respalda con 42 pruebas automatizadas, entre
ellas pruebas aleatorias que comparan cada operación con un modelo de
referencia independiente, y se comprobó que la suite detecta un defecto
introducido a propósito. Para la entrega final quedan la reclasificación
de urgencias, la carga por lotes, las estadísticas de atención y un
análisis experimental más amplio.
